\section{Analytic estimates for the square}\label{sec:moments}

It remains to prove the moment identities and detector bound in
Proposition~\ref{prop:weight-moments}. The products in $Z^2$ contain two
divisor-sum factors, those in $Z^4$ contain four, and the detector estimate
uses six. We first evaluate all these products in one formula that allows
repeated shifts and one prime mark. We then average the resulting singular
series over the shifts in $I$ and $J$, and apply the formulas to the square.
All dimensions, support budgets and functions remain fixed as $X$ tends
to infinity.

The leading coefficient separates into a test-function constant, determined
by which coordinates share a shift, and a singular series carrying the
numerical shifts. This separation is what permits the subset expansions
of $Z$ to be summed uniformly.

For a finite set $\mathcal H$ of distinct integers, put
\[
 \nu_p(\mathcal H)=|\mathcal H\bmod p|,\qquad
 \mathfrak S(\mathcal H)=\prod_p
 (1-\nu_p(\mathcal H)/p)(1-1/p)^{-|\mathcal H|}.
\]
We set $\mathfrak S(\varnothing)=1$. For a nonempty set the product converges:
beyond its diameter, the factors are $1+O_{|\mathcal H|}(p^{-2})$. It is nonnegative
and may be zero. The product is one when $|\mathcal H|=1$.

\begin{proposition}[Uniform mixed moments]\label{prop:mixed-moments}
Put $L=\log X$, where $X$ tends to infinity through integers, and let
$\mathbb E_X$ denote the average over $X<m\le2X$. Fix nonnegative integers
$j_1,\ldots,j_R$. For each $r$, let $F_r$ be a real function on
$[0,\infty)^{j_r}$ admitting a smooth compactly supported extension to
$\mathbb R^{j_r}$, and suppose that on the orthant it vanishes outside
\[
 v_1+\cdots+v_{j_r}\le\rho_r.
\]
Here the support budgets $\rho_r\ge0$ are fixed. Dimension zero means a scalar,
with support budget zero. For positive divisors
define
\[
 D_{F_r}(m;\mathbf b^{(r)})=
 \sum_{d_{r,\ell}\mid m+b_{r,\ell}\ (1\le\ell\le j_r)}
 \prod_{\ell=1}^{j_r}\mu(d_{r,\ell})
 F_r\left(\left(\frac{\log d_{r,\ell}}L\right)_{\ell=1}^{j_r}\right).
\]
For $j_r=0$, this means $D_{F_r}=F_r$.

Let $\Gamma=\{(r,\ell):1\le r\le R,\ 1\le\ell\le j_r\}$, and fix a
surjection $\kappa:\Gamma\to\{1,\ldots,t\}$. Write
$\mathcal M_i=\kappa^{-1}(\{i\})$. Take distinct integer shifts
$a_1,\ldots,a_t$ with $|a_i|\le DL^B$, where $D>0$ and $B\ge1$ are fixed,
and set $b_{r,\ell}=a_{\kappa(r,\ell)}$. Choose $\delta\in\{0,1\}$.
In the marked case take one more
shift $a_0$, distinct from these and satisfying the same bound. Put
\[
 (\chi_\delta(m),\mathcal H)=
 \begin{cases}
 (1,\{a_1,\ldots,a_t\}),&\delta=0,\\
 (\vartheta(m+a_0),\{a_0,a_1,\ldots,a_t\}),&\delta=1,
 \end{cases}
 \qquad \vartheta(n)=(\log n)1_{\{n\text{ prime}\}}.
\]
Assume $\sum_r\rho_r<1$ if $\delta=0$, and $\sum_r\rho_r<1/2$ if
$\delta=1$. Then there is a real constant $\mathcal C$, depending only on
the functions and $\kappa$, such that
\begin{equation}\label{eq:mixed-moment}
 \mathbb E_X\left[\chi_\delta(m)\prod_rD_{F_r}(m;\mathbf b^{(r)})\right]
 =L^{-t}\left\{\mathfrak S(\mathcal H)\mathcal C
 +O\!\left(L^{-1/2}(\log L)^A\right)\right\}.
\end{equation}
Here $A$ is a fixed finite number depending only on $\sum_rj_r$; the implied
constant may depend on all the fixed data, but is uniform in the actual
shifts. In particular the error is additive when $\mathfrak S(\mathcal H)=0$.

If each $|\mathcal M_i|$ is at most two, the constant is
\begin{equation}\label{eq:mixed-derivative-constant}
 \mathcal C=
 \int\prod_r\left((-1)^{j_r}\partial_1\cdots\partial_{j_r}F_r\right)
       \big((y_{\kappa(r,\ell)})_{\ell=1}^{j_r}\big)
       \prod_{i:\,|\mathcal M_i|=2}dy_i,
\end{equation}
where $y_i=0$ for $|\mathcal M_i|=1$, and each integrated variable ranges over
$[0,\infty)$. The zero-dimensional derivative is the original scalar.
When $\Gamma$ is empty, take $t=0$ and $\mathcal C=\prod_rF_r$.
\end{proposition}

The integer $t$ counts distinct numerical shifts, whereas
$\sum_r j_r$ counts divisor coordinates with repetitions. The power
$L^{-t}$ is determined by the former. Repetitions are recorded by the
fibers $\mathcal M_i$ and affect the constant $\mathcal C$. For two
divisor-sum factors whose individual shift tuples have distinct entries,
these fibers have size at most two. Products of four such factors
may have larger fibers; their estimates will use the general
finiteness assertion rather than the displayed derivative formula.
The marked formula retains the same power $L^{-t}$, consistent with
$\vartheta(m+a_0)$ having average tending to one. The additional factor
$L^{-1}$ in $V_B$ is supplied when the marked shifts are summed.

The scalar divisor-sum correlations of Goldston--Y{\i}ld{\i}r{\i}m
\cite{GoldstonYildirim2003Triple,GY2005} and Tao's smooth Fourier formulation,
described by Maynard \cite[Section~6, p.~404]{Maynard2015} motivate this calculation. We include its
proof because coupled smooth functions, arbitrary equality patterns and
an additive error uniform in logarithmic shifts are needed here.
The choice of smooth cutoffs was also informed by the analysis of
smoothing and high moments of divisor sums by
Granville--Koukoulopoulos--Maynard \cite{GKM2021}.  This is motivation
for the cutoff choice, not a source of the mixed-moment proposition
proved here.

\subsection{Residue classes and the prime mark}
Our objective is to replace a product of divisor sums, possibly weighted by one shifted prime, by an explicit finite density sum. We use the ordinary endpoint Bombieri--Vinogradov theorem in the form recorded by Bombieri, Friedlander and Iwaniec \cite[p.~206, equation~(1.4)]{BFI1986}; their $\psi$ notation is defined on p.~204. The additional divisor weights and endpoint shifts required here are handled below.

We retain the notation of Proposition~\ref{prop:mixed-moments} and put
$H_X=DL^B$. The two arithmetic lemmas below also allow $B\ge0$.
Negative shifts cause no difficulty, since all shifted arguments are
positive for sufficiently large $X$.
We write $\Lambda$ for the von Mangoldt function, $\varphi(q)$ for Euler's
totient, $\omega(q)$ for the number of distinct prime divisors of $q$,
$d(q)$ for the number of positive divisors, and $\zeta$ for the Riemann
zeta function.

\begin{lemma}[Weighted distribution with uniform shifted endpoints]\label{lem:arithmetic-1}

Assume the Bombieri--Vinogradov theorem: for every $A>0$ there is $b(A)>0$ such that
\begin{equation}
\sum_{q\le x^{1/2}(\log x)^{-b(A)}}
 \max_{(c,q)=1}\left|\psi(x;q,c)-\frac{x}{\varphi(q)}\right|
 \ll_A x(\log x)^{-A}.
\label{eq:arithmetic-BV}
\end{equation}
Here $\psi(x;q,c)=\sum_{n\le x,\ n\equiv c\pmod q}\Lambda(n)$, and $q=1$ is included. Fix $0\le\sigma<1/2$, $K\ge1$, and define
\[
E_X(q)=\max_{\substack{a\in\mathbb Z,\ |a|\le H_X\\(c,q)=1}}
\left|\sum_{\substack{X+a<n\le2X+a\\n\equiv c\pmod q}}\vartheta(n)
       -\frac{X}{\varphi(q)}\right|.
\]
For every $A>0$,
\begin{equation}
\sum_{q\le X^\sigma}\mu^2(q)K^{\omega(q)}E_X(q)
 \ll_{A,\sigma,K,D,B}X L^{-A}.
\label{eq:arithmetic-1}
\end{equation}
The conventions are $\omega(1)=0$, $K^0=1$, and $\varphi(1)=1$.

\end{lemma}
\begin{proof} At either endpoint $x=X,2X$, the range $q\le X^\sigma$ is contained in the range of \eqref{eq:arithmetic-BV} for every fixed logarithmic saving once $X$ is sufficiently large. Removing proper prime powers from a $\psi$-sum costs at most
\[
\sum_{p^j\le3X,\ j\ge2}\log p\ll X^{1/2}L^2
\]
in each residue class. Indeed there are $O(X^{1/2}L)$ pairs $p,j\ge2$, and each summand is $O(L)$. Translating the two endpoints through any $|a|\le H_X$ costs $O((H_X+1)L)$, uniformly in the residue class. Thus, writing $Q=\lfloor X^\sigma\rfloor$, for every $A_0>0$,
\begin{equation}
\sum_{q\le Q}E_X(q)
 \ll_{A_0}X L^{-A_0}
       +Q\{X^{1/2}L^2+(H_X+1)L\}
 \ll_{A_0,\sigma,D,B}X L^{-A_0}.
\label{eq:arithmetic-2}
\end{equation}
The last absorption uses the fixed strict inequality $\sigma<1/2$. The same bound controls the maximum over all shifts at once; no union bound over shifts is used.

An interval of length $X$ contains at most $X/q+1$ integers in one residue class. Also
\[
\frac q{\varphi(q)}=\prod_{p\mid q}\frac p{p-1}
 \le2^{\omega(q)}\le d(q).
\]
Since $q\le Q<X$ for large $X$, these two facts give
\begin{equation}
E_X(q)\ll \frac Xq\{L+d(q)\}.
\label{eq:arithmetic-3}
\end{equation}
For every fixed real $T\ge1$,
\begin{equation}
\sum_{q\le Q}\frac{\mu^2(q)T^{\omega(q)}}q
 \le\prod_{p\le Q}(1+T/p)\ll_T L^{c_T}
\label{eq:arithmetic-4}
\end{equation}
for a finite constant $c_T$. The elementary bound $\sum_{p\le y}1/p=O(\log\log(3y))$ suffices: for $y\ge2$, compare $1/p\le e p^{-1-1/\log y}$ with $\log\zeta(1+1/\log y)$, and bound $\zeta(1+u)\le1+1/u$. No prime-number theorem is needed for this estimate.

As squarefree $q$ satisfy $d(q)=2^{\omega(q)}$, \eqref{eq:arithmetic-3}--\eqref{eq:arithmetic-4} imply, for some finite $c_K$,
\[
\sum_{q\le Q}\mu^2(q)K^{2\omega(q)}E_X(q)
 \ll_K X L^{c_K}.
\]
Cauchy--Schwarz, applied to the nonnegative numbers $E_X(q)$, therefore gives
\[
\begin{split}
\sum_{q\le Q}\mu^2(q)K^{\omega(q)}E_X(q)
&\le
 \left(\sum_{q\le Q}E_X(q)\right)^{1/2}
 \left(\sum_{q\le Q}\mu^2(q)K^{2\omega(q)}E_X(q)\right)^{1/2}\\
&\ll X L^{(c_K-A_0)/2}.
\end{split}
\]
Choose $A_0\ge c_K+2A$. This proves \eqref{eq:arithmetic-1}. \end{proof}

\begin{lemma}[Exact density sum and arithmetic error]\label{lem:arithmetic-2}

Use the notation and support assumptions of
Proposition~\ref{prop:mixed-moments}; this lemma requires only bounded
functions $F_r$ and also allows $B\ge0$ in the shift bound
$H_X=DL^B$. Put $M=|\Gamma|$ and $\sigma=\sum_r\rho_r$.
As in the proposition, assume $\sigma<1$ when $\delta=0$ and
$\sigma<1/2$ when $\delta=1$.

For a squarefree divisor tuple $\mathbf d=(d_\gamma)_{\gamma\in\Gamma}$,
put $q=[d_\gamma:\gamma\in\Gamma]$, with $q=1$ when $M=0$.
Call the tuple compatible if, for every $p\mid q$, the shifts
$b_\gamma$ with $p\mid d_\gamma$ have one common residue modulo $p$.
When $\delta=1$, require this residue to differ from $a_0\bmod p$.
Set
\[
 \rho_\delta(\mathbf d)=
 \begin{cases}
  \displaystyle\prod_{p\mid q}(p-\delta)^{-1},
       &\mathbf d\text{ compatible},\\
  0,&\text{otherwise},
 \end{cases}
\]
and define the finite density sum
\begin{equation}
 \mathcal A_\delta=
 \sum_{\mathbf d\text{ squarefree}}
 \rho_\delta(\mathbf d)\prod_{\gamma\in\Gamma}\mu(d_\gamma)
 \prod_rF_r\left(
       \left(\frac{\log d_{r,\ell}}L\right)_{\ell=1}^{j_r}\right).
\label{eq:arithmetic-5}
\end{equation}
Empty products are one. For every $A>0$, uniformly in the allowed shifts,
\begin{equation}
 \mathbb E_X\left[\chi_\delta(m)
             \prod_rD_{F_r}(m;\mathbf b^{(r)})\right]
 =\mathcal A_\delta+O_A(L^{-A}).
\label{eq:arithmetic-6}
\end{equation}
The implied constant may depend on the fixed functions, dimensions,
support budgets, $D,B,A$. For $\delta=0$ and $M\ge1$, the error is
\begin{equation}
 O\bigl(X^{\sigma-1}(1+\log X)^{M-1}
          \prod_r\|F_r\|_\infty\bigr).
\label{eq:arithmetic-7}
\end{equation}

\end{lemma}
\begin{proof} Expand the product of divisor sums. M\"obius factors kill all nonsquarefree divisors. Whenever the product of the $F_r$'s is nonzero, the support hypothesis gives
\begin{equation}
\prod_{\ell=1}^{j_r}d_{r,\ell}\le X^{\rho_r}
\quad\text{for every }r,
\qquad
q\le\prod_\gamma d_\gamma\le X^\sigma.
\label{eq:arithmetic-8}
\end{equation}
This uses the budget of the sum of the coordinates in each factor, and does not impose a budget on the dimension.

For a squarefree tuple, the conditions $d_\gamma\mid m+b_\gamma$ prescribe at prime $p\mid q$ the residues $m\equiv-b_\gamma\pmod p$ for all selected coordinates. These congruences are inconsistent unless the selected shifts agree modulo every $p\mid q$. When they agree, the Chinese remainder theorem prescribes exactly one class $m\pmod q$. In the unmarked case this is precisely compatibility, and the count is $X/q+O(1)$, with an absolute constant in the $O(1)$.

When $M\ge1$, the number of positive integer tuples whose product is at most $Y\ge1$ is at most
\begin{equation}
Y\left(\sum_{d\le Y}\frac1d\right)^{M-1}
 \le Y(1+\log Y)^{M-1}.
\label{eq:arithmetic-9}
\end{equation}
Indeed, fix the first $M-1$ entries and bound the number of possible last entries by $Y$ divided by their product, then enlarge each preceding summation to $d\le Y$. Apply \eqref{eq:arithmetic-9} with $Y=X^\sigma$, multiply the $O(1)$ count error by $\prod_r\|F_r\|_\infty$, and divide by $X$. This proves \eqref{eq:arithmetic-7}, hence \eqref{eq:arithmetic-6} in the unmarked case. If $M=0$, the unmarked average and \eqref{eq:arithmetic-5} are both exactly $\prod_rF_r$.

For the marked case, put $n=m+a_0$. On a residue class satisfying these divisor congruences, the residue of $n\pmod p$ is $a_0-b_\gamma$ whenever $p\mid d_\gamma$. Consequently this class is reduced modulo $q$ exactly when the additional marked compatibility condition holds. If it is not reduced, a prime $n$ in that class would have to be a prime divisor of $q$, hence satisfy $n\le q$. But \eqref{eq:arithmetic-8} and $\sigma<1/2$ give
\[
n\ge X-H_X> X^\sigma\ge q
\]
for large $X$. The marked sum on a nonreduced class is therefore exactly zero, not an unestimated exceptional term. On a reduced class, the marked sum is
\[
\frac X{\varphi(q)}+O(E_X(q)).
\]
Since $q$ is squarefree, $\varphi(q)=\prod_{p\mid q}(p-1)$, proving the main density factor in \eqref{eq:arithmetic-5}.

For fixed squarefree $q$, each prime dividing $q$ must be assigned to a nonempty subset of the $M$ divisor coordinates. There are at most $(2^M-1)^{\omega(q)}\le (2^M)^{\omega(q)}$ choices if $M\ge1$; support and compatibility only reduce this count. The empty tuple, if $M=0$, has $q=1$ and is handled directly by the same bound with $K=1$. Thus the total normalized error is at most
\[
\frac{\prod_r\|F_r\|_\infty}{X}
 \sum_{q\le X^\sigma}\mu^2(q)(2^M)^{\omega(q)}E_X(q),
\]
which is $O_A(L^{-A})$ by Lemma~\ref{lem:arithmetic-1}. This proves the marked case, including all scalar factors and the wholly zero-dimensional case. \end{proof}

\subsection{Fourier evaluation of the density sum}

The residue-class calculation has controlled the prime-counting error.
It remains to evaluate its finite density sum, separating a constant
determined by the test functions from the singular series of the shifts.

\begin{proof}[Proof of Proposition~\ref{prop:mixed-moments}]
Let $M=|\Gamma|$. By Lemma~\ref{lem:arithmetic-2}, it remains to
evaluate $\mathcal A_\delta$ from \eqref{eq:arithmetic-5}: the arithmetic
error is $O_N(L^{-N})$ for every fixed $N$.
If $M=0$, that sum is exactly $\prod_rF_r$, and the assertion follows.
We henceforth assume $M>0$.

Choose a smooth compact extension $\widetilde F_r$ of each positive-dimensional
function and define its Fourier transform by
\[
 \Phi_r(\mathbf u)=(2\pi)^{-j_r}\int_{\mathbb R^{j_r}}
 e^{\sum_\ell v_\ell}\widetilde F_r(\mathbf v)
 e^{\mathrm i\mathbf u\cdot\mathbf v}\,d\mathbf v.
\]
The function $\Phi_r$ is Schwartz, and Fourier inversion gives exactly
\begin{equation}\label{eq:fourier-inversion}
 F_r(\mathbf v)=\int_{\mathbb R^{j_r}}\Phi_r(\mathbf u)
 e^{-\sum_\ell(1+\mathrm i u_\ell)v_\ell}\,d\mathbf u
 \qquad(\mathbf v\in[0,\infty)^{j_r}).
\end{equation}
For scalar factors use the scalar itself. Put $\Phi=\prod_r\Phi_r$,
$z_\gamma=(1+\mathrm i u_\gamma)/L$, and
$z_S=\sum_{\gamma\in S}z_\gamma$.

Insert \eqref{eq:fourier-inversion} into the density sum
\eqref{eq:arithmetic-5}, retaining the sum over all squarefree tuples.
The following bound justifies moving that sum inside the integrals:
\begin{equation}\label{eq:absolute-euler-sum}
 \sum_{\mathbf d}\rho_\delta(\mathbf d)\prod_\gamma d_\gamma^{-1/L}
 \le\prod_p(1+C_Mp^{-1-1/L})\ll_M L^{C_M}.
\end{equation}
Indeed $1/(p-\delta)\le2/p$, every nonempty selected subset contributes at
most $p^{-1/L}$, and there are at most $2^M-1$ subsets. Moreover
$\sum_pp^{-1-\varepsilon}\le\log\zeta(1+\varepsilon)
\le\log(1+1/\varepsilon)$. Multiplication by the finite integral
$\int|\Phi|$ proves absolute convergence of the sum-integral interchange.

Call a nonempty subset $S\subset\Gamma$ allowed at $p$ if its shifts satisfy
the same compatibility and mark-exclusion conditions. The exact Euler factor
is then
\[
 P_p(\mathbf z)=1+\frac1{p-\delta}
 \sum_{S\text{ allowed at }p}(-1)^{|S|}p^{-z_S},
 \qquad
 \mathcal A_\delta=\int\Phi(\mathbf u)\prod_pP_p(\mathbf z)\,d\mathbf u.
\]
By \eqref{eq:absolute-euler-sum} and Schwartz decay, the part of this integral
where $\max_\gamma|u_\gamma|>\sqrt L$ is $O_N(L^{-N})$ for every fixed $N$.
All subsequent comparisons are made on the complementary region.

The factors carrying the zeta poles are
\[
 B_p(\mathbf z)=\prod_{i=1}^t\prod_{\varnothing\ne S\subset\mathcal M_i}
 (1-p^{-1-z_S})^{-(-1)^{|S|}},\qquad
 H_p(\mathbf z)=P_p(\mathbf z)/B_p(\mathbf z).
\]
For $\Re z_\gamma\ge0$, all denominators in this expression have absolute
value at least $1/2$, and $|H_p(\mathbf z)|\le\exp(C_M/p)$. If distinct shifts
in $\mathcal H$ have distinct residues modulo $p$, the allowed subsets are
exactly the nonempty subsets of the individual fibers $\mathcal M_i$.
Matching the terms of order $p^{-1}$ gives
\begin{equation}\label{eq:generic-local-factor}
 H_p(\mathbf z)=1+O_M(p^{-2}),
\end{equation}
uniformly for $\Re z_\gamma\ge0$.

At $\mathbf z=0$, summing signs over the nonempty subsets of each occupied
allowed residue class gives
\[
 P_p(0)=\frac{p-\nu_p(\mathcal H)}{p-\delta},\qquad
 B_p(0)=(1-1/p)^t.
\]
It follows that $H_p(0)$ is the local factor of $\mathfrak S(\mathcal H)$.
This identity holds also when that factor is zero.

We now compare the product of the $H_p$'s additively with its value at zero.
Take $Y=\max(2,2DL^B)$. For $p\le Y$, differentiation along the segment from
zero to $\mathbf z$ gives
\[
 |H_p(\mathbf z)-H_p(0)|\ll_M L^{-1/2}\frac{\log p}{p}.
\]
For example, $|p^{-z_S}-1|\le|z_S|\log p$ on that segment, while the denominators
stay bounded away from zero. Each mixed partial product in the additive
telescoping identity is at most
\[
 \exp\left(C_M\sum_{p\le Y}1/p\right)\ll_M(\log Y)^{C'_M}.
\]
Using even the elementary bound
$\sum_{p\le Y}(\log p)/p\le\sum_{n\le Y}(\log n)/n=O((\log Y)^2)$, we obtain
an error $O(L^{-1/2}(\log L)^A)$ for the finite product. For $p>Y$ there are
no collisions of distinct shifts, so \eqref{eq:generic-local-factor} shows
that both remaining tail products are $1+O_M(1/Y)$. Consequently
\begin{equation}\label{eq:zero-safe-comparison}
 \prod_pH_p(\mathbf z)=\mathfrak S(\mathcal H)
  +O\!\left(L^{-1/2}(\log L)^A\right),\qquad
 \mathfrak S(\mathcal H)\ll(\log L)^A.
\end{equation}
This argument never divides by a possibly zero local singular-series factor.

It remains to evaluate the zeta factors. Absolute Euler products for
$\Re z_\gamma>0$ give
\[
 \prod_pB_p(\mathbf z)=\prod_{i=1}^t
 \prod_{\varnothing\ne S\subset\mathcal M_i}
 \zeta(1+z_S)^{(-1)^{|S|}}.
\]
Write $w_\gamma=1+\mathrm i u_\gamma$ and
\[
 K_i(\mathbf u)=\prod_{\varnothing\ne S\subset\mathcal M_i}
 \left(\sum_{\gamma\in S}w_\gamma\right)^{(-1)^{|S|+1}}.
\]
The pole expansion $\zeta(1+z)=z^{-1}(1+O(|z|))$ and the identity
$\sum_{\varnothing\ne S\subset\mathcal M_i}(-1)^{|S|}=-1$ imply
\begin{equation}\label{eq:zeta-kernel}
 \prod_pB_p(\mathbf z)
 =L^{-t}\prod_iK_i(\mathbf u)\{1+O_M(L^{-1/2})\}.
\end{equation}
Every denominator in $K_i$ has positive integer real part, so the product of
the $K_i$ has at most polynomial growth. Thus
\[
 \mathcal C=\int_{\mathbb R^M}\Phi(\mathbf u)\prod_iK_i(\mathbf u)\,d\mathbf u
\]
is absolutely convergent. Integrating \eqref{eq:zero-safe-comparison} and
\eqref{eq:zeta-kernel} against the Schwartz transform, then restoring the
negligible frequency tails, proves \eqref{eq:mixed-moment}. The constant is
real by conjugation under $\mathbf u\mapsto-\mathbf u$.

Finally, a singleton fiber contributes $w_\gamma$, and a double fiber
contributes
\[
 \frac{w_\gamma w_{\gamma'}}{w_\gamma+w_{\gamma'}}
 =w_\gamma w_{\gamma'}\int_0^\infty
 e^{-(w_\gamma+w_{\gamma'})y_i}\,dy_i.
\]
After taking absolute values, the new variables contribute
$e^{-2\sum_i y_i}$; the remaining factor is a Schwartz function times a fixed
polynomial. Fubini therefore applies. Recombining the Fourier integrals by
differentiating \eqref{eq:fourier-inversion} proves
\eqref{eq:mixed-derivative-constant}.
\end{proof}

\subsection{Averaging the singular series}\label{sec:singular-average}

We next average the numerical-shift factor in
Proposition~\ref{prop:mixed-moments}. The following form allows different
coordinates to range over either of the two blocks. It is a box version
of Gallagher's singular-series mean \cite[p.~5, equation~(3)]{Gallagher1976}.
We give a self-contained proof by truncating the Euler product;
no statement about the distribution of primes in individual short intervals
is assumed.

\begin{lemma}[Singular series in boxes]\label{lem:singular-average}
Fix an integer $s\ge0$ and a constant $K\ge1$. For each positive integer $h$,
let $I_1,\ldots,I_s$ be intervals of $h$ consecutive integers, all contained
in an interval of diameter at most $Kh$. Then, as $h\to\infty$,
\[
 \sum_{\substack{a_i\in I_i\ (1\le i\le s)\\a_1,\ldots,a_s\text{ distinct}}}
 \mathfrak S(\{a_1,\ldots,a_s\})=h^s(1+o(1)).
\]
The error is uniform in the intervals subject to this condition, with $s,K$
fixed. Here $A_s$ denotes a fixed finite constant depending only on $s$. In fact the normalized sum differs from one by
\[
 O_{s,K}(1/\log\log h)+O_s\big(h^{-1/2}(\log\log h)^{A_s}\big).
\]
\end{lemma}

\begin{proof}
The cases $s=0,1$ are immediate. Otherwise let $y=(1/4)\log h$ and
$Q=\prod_{p\le y}p$. The prime number theorem gives $Q\le h^{1/2}$ for large
$h$. For all tuples, including those with equal entries, define
\[
 \mathfrak S_y(\mathbf a)=\prod_{p\le y}
 (1-\nu_p(\{a_1,\ldots,a_s\})/p)(1-1/p)^{-s}.
\]
The exponent remains $-s$ even on tuples with collisions. For independent
uniform residues $a_1,\ldots,a_s$ modulo a prime $p$, averaging over an
additional uniform residue $x$ gives
\[
 \mathbb E(1-\nu_p/p)
 =\mathbb P(a_i\ne x\text{ for every }i)=(1-1/p)^s.
\]
The Chinese remainder theorem therefore shows that $\mathfrak S_y$ has
exactly mean one on complete residue tuples modulo $Q$.

A uniform integer in a length-$h$ interval has residue distribution modulo
$Q$ within total variation $O(Q/h)$ of uniform. The elementary bound $\sum_{p\le y}1/p=O(\log\log(3y))$ used above gives
\[
 0\le\mathfrak S_y(\mathbf a)
 \le\prod_{p\le y}(1-1/p)^{-s}\ll_s(\log y)^{A_s}.
\]
It follows that
\[
 h^{-s}\sum_{a_i\in I_i}\mathfrak S_y(\mathbf a)
 =1+O_s\big((Q/h)(\log y)^{A_s}\big).
\]
There are $O_s(h^{s-1})$ tuples with equal entries, so deleting them changes
this normalized mean by at most $O_s(h^{-1}(\log y)^{A_s})$.

For a distinct tuple put $\Delta=\prod_{i<j}|a_i-a_j|$, a positive integer.
For $y>2s$ every factor of the tail
\[
 T_y(\mathbf a)=\prod_{p>y}(1-\nu_p/p)(1-1/p)^{-s}
\]
is positive. A prime not dividing $\Delta$ has $\nu_p=s$ and a local factor
$1+O_s(p^{-2})$; the logarithm of a remaining factor is $O_s(1/p)$. Hence
\[
 |\log T_y|\ll_s \frac1y+\sum_{\substack{p>y\\p\mid\Delta}}\frac1p
 \le C_s\left(\frac1y+\frac{\log\Delta}{y\log y}\right).
\]
Since $\log\Delta\le\binom s2\log(Kh)$, this gives, uniformly in the distinct
tuples,
\[
 T_y(\mathbf a)=1+O_{s,K}(1/\log\log h).
\]
Now use the identity $\mathfrak S(\{a_1,\ldots,a_s\})=\mathfrak S_y(\mathbf a)T_y(\mathbf a)$
and the nonnegative truncated mean. This identity is valid also if a
small-prime factor vanishes; no quotient by $\mathfrak S_y$ is taken.
The stated error follows from $Q\le h^{1/2}$ and $\log y\asymp\log\log h$.
\end{proof}

When the shifts are summed, an equality pattern with $t$ distinct unmarked shifts has
$O(h^t)$ assignments in fixed blocks of length $h\asymp L$. Thus the uniform
error in Proposition~\ref{prop:mixed-moments}, summed over that pattern,
is $O(L^{-1/2}(\log L)^A)=o(1)$. If one additional distinct prime shift is
summed and its weight divided by $L$, there are $O(h^{t+1})$ assignments and
the same conclusion holds. Only finitely many patterns occur for a fixed
family of functions. The family is chosen before $X$ tends to infinity.

\subsection{Evaluation of the square}\label{sec:weight-moment-proof}

We now apply the correlation formula and singular-series mean to the
functions constructed in Section~\ref{sec:weights}. The signed complete
mixed derivative of each cumulative function is $f_j$, as in
\eqref{eq:cumulative-derivative}; its vanishing on coordinate faces will
leave only matching subsets in the quadratic moments.

The support restrictions in the moment formula apply to the sum of
the support budgets of the factors, independently of their dimensions.
For the present choice $\tau=1/8$, the required uses are as follows:
\begin{center}
\begin{tabular}{lll}
Quantity & Moment used & Total support budget\\
\hline
$\mathbb E_XZ^2$ & unmarked & $2\tau<1$\\
$\mathbb E_X(Z^2V_B)$, $B=I,J$ & one prime mark & $2\tau<1/2$\\
$\mathbb E_XZ^4$ & unmarked & $4\tau<1$\\
Detector off-diagonal terms & unmarked & $6\tau<1$
\end{tabular}
\end{center}
The last line uses divisor-sum majorants for the two possible primes
in $J$. It therefore requires no asymptotic formula with two prime marks.

\begin{proof}[Proof of Proposition~\ref{prop:weight-moments}]
First we check the smoothness and support needed by the moment
formula.  Extend $f_j$ by zero to $\mathbb R^j$.  Its cumulative integral
is smooth on $\mathbb R^j$, and multiplication by smooth coordinate
cutoffs that equal one near $[0,\infty)$ and vanish below $-1$ makes it
compactly supported without changing its orthant values or derivatives.
Indeed, the cumulative integral already vanishes when any coordinate
exceeds an upper support bound of $f_j$.  On the nonnegative orthant,
$F_j$ vanishes if $\sum_i v_i\ge\tau$, and its signed complete mixed derivative is given by
\eqref{eq:cumulative-derivative}.
This derivative is zero on every coordinate face.  For $j\ge1$,
$Tf_{j+1}$ has the same smoothness, interior-support, and face-vanishing
properties.

Expand $Z^2$ over pairs of subsets $S,T\subset I$.  Any shift in
$S\mathbin\triangle T$ occurs once and sets a coordinate of one of the
functions in \eqref{eq:cumulative-derivative} to zero.  Its main
coefficient is therefore zero.  If $S=T$ has size $j$, the coefficient
is $\|f_j\|_2^2$.  Averaging the singular series and using
\[
 \frac{\binom hj}{L^j}\longrightarrow\alpha_j
\]
proves the first limit in \eqref{eq:weight-unmarked}.  With a prime mark
in $J$, the same matching rule applies and
\[
 \frac{h\binom hj}{L^{j+1}}\longrightarrow\lambda\alpha_j.
\]
The two divisor factors have total support $2\tau<1/2$, so the marked
moment formula proves the first limit in \eqref{eq:weight-marked}.

For an internal mark $a\in I$, the exact deletion identity
\eqref{eq:prime-deletion} replaces $F_j$ by
$\widetilde F_j=F_j+F_{j+1}(\cdot,0)$ on the prime event.
Deletion is a bijection, so no extra dimension factor occurs.  The
support radius of $\widetilde F_j$ is still $\tau$, and its signed
complete mixed derivative is $f_j+Tf_{j+1}$.  It vanishes on coordinate
faces when $j\ge1$.  Apply the external marked formula to these
functions.  Only equal subsets $U$ survive, and
\[
 \frac{h\binom{h-1}{j}}{L^{j+1}}\longrightarrow\lambda\alpha_j.
\]
This proves the formula for $v$, including the scalar term
$\lambda|f_0+\int f_1|^2$.

The equality-pattern estimate preceding this subsection justifies
summing the errors in these limits. Applied to $Z^4$, the same argument
uses four divisor factors, total support
$4\tau<1$, and only finiteness of their pattern constants.  Bounded
averages of the singular series then give the second assertion in
\eqref{eq:weight-unmarked}.

It remains to prove a detector estimate whose constant is independent
of the chosen family.  Write $q_X=\log(2X+2h)/L=1+o(1)$.  At a
relevant prime $m+b$, we have $D_G(m;b)=1$.  Consequently, pointwise,
\[
 \frac{\vartheta(m+b)}L\le q_XD_G(m;b)^2,
 \qquad
 \left(\frac{\vartheta(m+b)}L\right)^2
       \le q_X\frac{\vartheta(m+b)}L.
\]
The first inequality holds at composites as well, since its left side
is zero.  Because $Z^2\ge0$, these inequalities give
\begin{align*}
 \mathbb E_X(Z^2V_J^2)
 &\le q_X\mathbb E_X(Z^2V_J)\\
 &\quad+q_X^2\sum_{\substack{b,c\in J\\b\ne c}}
       \mathbb E_X\bigl[Z^2 D_G(m;b)^2D_G(m;c)^2\bigr].
\end{align*}
An off-diagonal summand contains six divisor factors, with total
support $6\tau<1$.  The shifts $b,c$ are paired and lie outside $I$.
The two subsets from $Z^2$ must again be identical.  For their common
size $j$, the coefficient is $c_G^2\|f_j\|_2^2$.  The pair $(b,c)$ is
ordered, and
\[
 \frac{h(h-1)\binom hj}{L^{j+2}}
      \longrightarrow\lambda^2\alpha_j.
\]
Thus the off-diagonal sum tends to $\lambda^2c_G^2w$.
Together with \eqref{eq:weight-marked}, this proves
\eqref{eq:detector-bound}.  Although the convergence rates may depend
on the fixed family, the resulting $C_*$ does not.
\end{proof}
